% uid: ""
% school: "lycee"
% level: "1ere"
% chapter: "Activités numériques"
% title: "Racines irrationnelles"
% difficulty: "Difficile"
% tags: ["racines-carrees", "irrationalite", "raisonnement-par-l-absurde", "parite"]
% provenance: "own"
% batch: "2026-11-22-revision-1ere-A"

% ---STATEMENT---
Un nombre est \emph{rationnel} s'il s'écrit $\dfrac pq$ avec $p$, $q$ entiers, $q\neq0$, la fraction étant irréductible ($p$ et $q$ n'ont pas de diviseur commun autre que $1$).
\begin{enumerate}
  \item Montrer que si le carré d'un entier $n$ est pair, alors $n$ est pair. [Étudier le cas $n$ impair.]
  \item On suppose que $\sqrt2=\dfrac pq$ (fraction irréductible). Montrer que $p^2=2q^2$, puis que $p$ et $q$ sont pairs. Conclure que $\sqrt2$ n'est pas rationnel.
  \item Montrer de même que $\sqrt6$ n'est pas rationnel.
  \item En calculant $(\sqrt2+\sqrt3)^2$, montrer que $\sqrt2+\sqrt3$ n'est pas rationnel.
  \item En déduire que $\sqrt3-\sqrt2$ n'est pas rationnel.
\end{enumerate}

% ---SOLUTION---
\begin{prop}[Raisonnement par l'absurde]
Pour montrer qu'une propriété est vraie, on suppose qu'elle est \textbf{fausse} et on aboutit à une \ib{contradiction}. Un entier est pair s'il s'écrit $2m$, impair s'il s'écrit $2m+1$ ; le produit de deux entiers impairs est impair.
\end{prop}
\begin{enumerate}
  \item On montre : \og $n$ impair $\implies n^2$ impair \fg{}.

  Si $n=2m+1$ :
  \[n^2=4m^2+4m+1=2\,(2m^2+2m)+1\]
  $n^2$ est impair. Donc, si $n^2$ est pair, $n$ ne peut pas être impair : \textbf{$n$ est pair}.

  \item On élève $\sqrt2=\dfrac pq$ au carré : $2=\dfrac{p^2}{q^2}$, soit
  \[p^2=2q^2\]
  Donc $p^2$ est pair, et d'après la question 1, $p$ est pair : $p=2m$.

  On remplace : $4m^2=2q^2$, donc
  \[q^2=2m^2\]
  Donc $q^2$ est pair, et $q$ est pair.

  $p$ et $q$ sont tous deux pairs : $2$ est un diviseur commun, ce qui contredit \og fraction irréductible \fg{}. L'hypothèse est fausse :
  \[\ib{$\sqrt2$ n'est pas rationnel}\]

  \item On suppose $\sqrt6=\dfrac pq$ (irréductible). Alors $p^2=6q^2$ : $p^2$ est pair, donc $p=2m$.

  On remplace : $4m^2=6q^2$, soit
  \[2m^2=3q^2\]
  Le nombre $3q^2$ est pair. Si $q^2$ était impair, $3q^2$ serait un produit de deux impairs, donc impair : c'est faux. Donc $q^2$ est pair, et $q$ est pair.

  $p$ et $q$ sont pairs : contradiction. \ib{$\sqrt6$ n'est pas rationnel.}

  \item On développe :
  \[(\sqrt2+\sqrt3)^2=2+2\sqrt6+3=5+2\sqrt6\]
  Supposons $\sqrt2+\sqrt3=r$ rationnel. Alors $r^2=5+2\sqrt6$, donc
  \[\sqrt6=\dfrac{r^2-5}2\]
  Ce nombre est rationnel (opérations sur des rationnels) : cela contredit la question 3.

  \ib{$\sqrt2+\sqrt3$ n'est pas rationnel.}

  \item On calcule le produit conjugué :
  \[(\sqrt3-\sqrt2)(\sqrt3+\sqrt2)=3-2=1\implies\sqrt3-\sqrt2=\dfrac1{\sqrt3+\sqrt2}\]
  Supposons $\sqrt3-\sqrt2=s$ rationnel. Comme $s>0$, $\sqrt3+\sqrt2=\dfrac1s$ serait rationnel : cela contredit la question 4.

  \ib{$\sqrt3-\sqrt2$ n'est pas rationnel.}
\end{enumerate}
\textbf{Idée :} supposer le contraire et chercher un diviseur commun caché dans la fraction \og irréductible \fg{}.

\textbf{Piège :} pour $\sqrt6$, on aboutit à $2m^2=3q^2$ : il faut encore justifier que $q$ est pair (le facteur $3$ est impair).
